\documentclass[10pt,a4paper]{amsart}
\usepackage[margin=19mm]{geometry}
\usepackage{amsmath,amssymb,mathtools}
\usepackage[scr=rsfs,cal=euler]{mathalfa}
\usepackage{xfrac}
\usepackage[unicode,hidelinks,bookmarks=false]{hyperref}
\newcommand{\ie}{i.e.\ }
\newcommand{\cf}{cf.\ }
\newcommand{\resp}{respectively\ }
\newcommand{\Subsection}{\subsection}
\providecommand{\nobreakdash}{\nobreak\hbox{-}}
\setlength{\emergencystretch}{2em}
\multlinegap0pt
\usepackage{mathtools}
\usepackage{amssymb}
\usepackage{enumerate}
\usepackage{breqn}
\newtheorem{theorem}{Theorem}
\newtheorem{corollary}{Corollary}
\title{Quadratic maps between non-abelian groups}
\author{Asgar Jamneshan, Andreas Thom}
\date{Source: arXiv:2412.14908v3 (CC BY 4.0)}
\begin{document}
\maketitle

\noindent\emph{Source and licence.} Quadratic maps between non-abelian groups. Version arXiv:2412.14908v3 (CC BY 4.0), distributed by arXiv under the Creative Commons Attribution 4.0 International licence, http://creativecommons.org/licenses/by/4.0/, as recorded in the arXiv licence metadata for this exact version and shown on the abstract page https://arxiv.org/abs/2412.14908v3. Full licence notice: \texttt{licenses/arxiv-2412.14908-CC-BY-4.0.txt}.\par\smallskip

\noindent\textit{Editorial scope.} Counted unit: the article's theorem thm:approx (source0.tex lines 226--229). The article's complete original proof of it is delivered with it (source0.tex lines 684--703, one \texttt{proof} environment of 20 lines, which the article places away from the statement). No mathematics is written, repaired or re-derived: the statement and the proof are the article's own lines, in the article's own notation, and no retained line is substituted or re-flowed.  Retained uncounted as the source-local context the proof needs: lines 669--673.  Nothing was omitted from any retained span, so the package carries no omission record.  No retained line contains a citation, so the wrapper carries no bibliography and no bibliography entry.  No retained line contains a figure environment, an image-inclusion command or drawing code, so no image is delivered and the package declares no asset dependency.  Editorial formatting only, in addition: an AMS \texttt{amsart} preamble that restates the article's own declarations and macros used by the retained lines, namely the environments \texttt{theorem}, \texttt{corollary} on separate counters, together with the standard packages those lines need.  The retained environments print in this document as Theorem 1, Corollary 1 in order of appearance, whereas the article prints the same items with the numbering of its own full document.  The complete native source of the article as distributed in the arXiv e-print for this exact version is bundled unchanged as \texttt{sources/170413/source0.tex}.
\begin{corollary} \label{cor:bound} Let $k \in \mathbb N.$
Let $G$ be a perfect group of commutator width $k$. Consider the free group on the set $\varphi(G)$. There exists a constant $c:=c(k)\geq 0$, such that for every $g,h \in G$, the element $\varphi(g_1g_2)\varphi(g_2)^{-1} \varphi(g_1)^{-1}$ is a product of $c$ conjugates of elements of the form
$$\varphi(g_1g_2g_3)^{-1} \varphi(g_1g_2)\varphi(g_2)^{-1}\varphi(g_1)^{-1} \varphi(g_1g_3)\varphi(g_3)^{-1}\varphi(g_2g_3)$$
or their inverses for $g_1,g_2,g_3 \in G$ .
\end{corollary}

\begin{theorem} \label{thm:approx}
Let $k \in \mathbb{N}$. There exists a constant $c := c(k) > 0$ such that the following holds:  
Let $d \in \mathbb{N}$, let $(H, \partial)$ be a group equipped with a bi-invariant metric, and let $G$ be a group of commutator width at most $k$. Then every unital uniform $\varepsilon$-polynomial map $\varphi \colon G \to H$ of degree~$d$ is a uniform $c(k)^{d-1} \varepsilon$-homomorphism.
\end{theorem}

\begin{proof}[Proof of Theorem \ref{thm:approx}]
The proof proceeds by induction on $d$, where the case $d=1$ is trivial. The case $d=2$ requires special attention. Let's compute:
\begin{eqnarray*}
(\Delta_{g_1,g_2,g_3} \varphi)(g_0) &=& (\Delta_{g_1}(\Delta_{g_2,g_3}  \varphi))(g_0) \\
&=& (\Delta_{g_2,g_3}  \varphi)(g_1g_0) ((\Delta_{g_2,g_3}  \varphi)(g_0))^{-1} \\
&=&\varphi(g_3g_2g_1g_0) \varphi(g_2g_1g_0)^{-1} \varphi(g_1g_0) \varphi(g_3g_1g_0)^{-1} \\
&& \cdot (\varphi(g_3g_2g_0) \varphi(g_2g_0)^{-1} \varphi(g_0) \varphi(g_3g_0)^{-1})^{-1} \\
&=&\varphi(g_3g_2g_1g_0) \varphi(g_2g_1g_0)^{-1} \varphi(g_1g_0) \varphi(g_3g_1g_0)^{-1} \\
&& \cdot \varphi(g_3g_0)\varphi(g_0)^{-1}\varphi(g_2g_0)\varphi(g_3g_2g_0)^{-1}
\end{eqnarray*}
Thus, for $\varphi$ unital, we obtain:
\begin{equation*}
(\Delta_{g_1,g_2,g_3} \varphi)(1)=
\varphi(g_3g_2g_1) \varphi(g_2g_1)^{-1} \varphi(g_1) \varphi(g_3g_1)^{-1} \varphi(g_3)\varphi(g_2)\varphi(g_3g_2)^{-1}.
\end{equation*}
By our assumption, we obtain
$$\partial(\varphi(g_3g_2g_1) \varphi(g_2g_1)^{-1} \varphi(g_1) \varphi(g_3g_1)^{-1} \varphi(g_3)\varphi(g_2)\varphi(g_3g_2)^{-1},1) \leq \varepsilon$$ for all $g_1,g_2,g_3 \in G$. Note that the same estimate holds for conjugates and inverses by bi-invariance of the metric. Thus, Corollary \ref{cor:bound} implies that $\partial(\varphi(g_1g_2)\varphi(g_2)^{-1}\varphi(g_1)^{-1},1)\leq c(k) \varepsilon$ and equivalently
$$\partial((\Delta_{g_1,g_2}\varphi)(1),1) \leq c(k) \varepsilon.$$
As required, we conclude that $\varphi$ is a uniform $c(k)\varepsilon$-polynomial of degree $1$. The general case follows by induction. Indeed, note that for $d \geq 2$, $\varphi$ is a uniform $\varepsilon$-polynomial of degree $d+1$ if and only if $\Delta_g \varphi$ is a uniform $\varepsilon$-polynomial of degree $d$ for all $g \in G.$
\end{proof}

\end{document}
